\documentclass{article}
\title{RRS Route Editor Math}
\author{Авилкин Николай}
\usepackage[none]{hyphenat}
\usepackage{indentfirst}
%\usepackage[utf8]{inputenc}
%\usepackage[T2A]{fontenc}
\usepackage[english,russian]{babel}
\usepackage{amsmath}
\usepackage{hyperref}
\begin{document}
	\maketitle
	\newpage
	\tableofcontents
	\newpage
	\section{Нормализация координат мыши}
	В VSG диапазон координат мыши:
	
	$x:[0;windowWidth]$
	
	$y:[0;windowHeight]$\\
	
	Чтобы сконвертировать их к нормализованным координатам $[-1.0;1.0]$, проводим
	следующие трансформации:
	\begin{equation}
		\begin{split}
			xNorm=\frac{x}{windowWidth}*2.0-1.0\\
			yNorm=\frac{y}{windowHeight}*2.0-1.0
		\end{split}
	\end{equation}
	\section{Мировые координаты мыши}
	Чтобы математически определить точку пересечения мыши с объектами сцены, нам
	необходимо получить луч, исходящий из мировых координат мыши вглубь сцены.
	
	Для этого нормализуем координаты мыши и добавляем к ним компонент $Z$ (глубину):
	обычно 0.0 для начала луча и 1.0 для точки, которая глубже/дальше относительно
	камеры и помогает определить направление луча.
	
	Потом последовательно умножаем полученные 3D-векторы (экранные координаты) на
	обратную матрицу проекции и на обратную матрицу вида. Получаем 2 точки в мировых
	координатах, которые как раз и образуют луч, который понадобится нам в
	дальнейших расчетах.
	\section{Параметрическое уравнение прямой}
	Для последующих расчётов будет удобно представлять нашу прямую (луч от мыши) в
	следующем виде:
	\begin{equation}
		\begin{cases}
			x=x_0+a_xt\\
			y=y_0+a_yt\\
			z=z_0+a_zt
		\end{cases}
	\end{equation}
	
	где $\vec{a}$ - это направляющий вектор, а $t$ - это вещественный параметр (при
	нормализованном векторе $\vec{a}$ это длина отрезка от начала луча до точки
	пересечения).
	\newpage
	\section{Пересечение прямой и плоскости}
	Мы будем рассчитывать пересечение с плоскостью, заданной точкой, лежащей на ней
	(обычно положением гизмо) и вектором нормали $\vec{n}$ (обычно направлением
	камеры).
	
	Уравнение плоскости, проходящей через точку $G(G_x;G_y;G_z)$ перпендикулярно
	вектору $\vec{n}(n_x;n_y;n_z)$, выражается формулой:
	\begin{equation}
		n_x(x-G_x)+n_y(y-G_y)+n_z(z-G_z)=0
	\end{equation}
	
	Назовём начало луча точкой $M(x_0;y_0;z_0)$.
	
	Подставим в это уравнение наши $x$, $y$ и $z$ из параметрического уравнения
	прямой, сократим и выразим $t$:
	\begin{equation}
		\begin{split}
			&n_x(x_0+a_xt-G_x)+n_y(y_0+a_yt-G_y)+n_z(z_0+a_zt-G_z)=0\\
			&n_xa_xt+n_ya_yt+n_za_zt=-n_xx_0+n_xG_x-n_yy_0+n_yG_y-n_zz_0+n_zG_z\\
			&t(\vec{n}\cdot\vec{a})=\vec{n}\cdot\vec{G}-\vec{n}\cdot\vec{M}\\
			&t=\frac{\vec{n}\cdot\vec{G}-\vec{n}\cdot\vec{M}}{\vec{n}\cdot\vec{a}}
		\end{split}
	\end{equation}
	
	Подставляем $t$ в параметрическое уравнение и получаем необходимую нам точку
	пересечения.
	\newpage
	\section{Пересечение прямой и цилиндра}
	Рассмотрим цилиндр, направленный вдоль оси $X$, центр основания которого
	расположен в точке $C(C_x;C_y;C_z)$. Радиус - $R$; высота - $H$. Он выражается
	формулами:
	\begin{equation}
		\begin{cases}
			(y-C_y)^2+(z-C_z)^2=R^2\\
			C_x\leq{}x\leq{}C_x+H
		\end{cases}
	\end{equation}
	
	Назовём начало луча точкой $M(x_0;y_0;z_0)$.
	
	Подставим в это уравнение наши $x$, $y$ и $z$ из параметрического уравнения
	прямой, доведём до квадратного уравнения:
	\begin{equation}
		\begin{split}
			&(y_0+a_yt-C_y)^2+(z_0+a_zt-C_z)^2=R^2\\\\
			&y_0^2+a_y^2t^2+C_y^2+2y_0a_yt-2y_0C_y-2a_ytC_y\\
			&+z_0^2+a_z^2t^2+C_z^2+2z_0a_zt-2z_0C_z-2a_ztC_z=R^2\\\\
			&t^2(a_y^2+a_z^2)+t(2y_0a_y-2a_yC_y+2z_0a_z-2a_zC_z)\\
			&+1(y_0^2+C_y^2-2y_0C_y+z_0^2+C_z^2-2z_0C_z-R^2)=0\\\\
			&t^2(a_y^2+a_z^2)+t(2(a_y(y_0-C_y)+a_z(z_0-C_z)))\\
			&+1((y_0-C_y)^2+(z_0-C_z)^2-R^2)=0
		\end{split}
	\end{equation}
	
	Обозначим $y_0-C_y$ и $z_0-C_z$ как $V_y$ и $V_z$ соответственно, тогда
	уравнение приобретёт следующий вид:
	\begin{equation}
		t^2(a_y^2+a_z^2)+t(2(a_yV_y+a_zV_z))+1(V_y^2+V_z^2-R^2)=0
	\end{equation}
	
	Решаем квадратное уравнение и находим t:
	\begin{equation}
		\begin{split}
			A&=a_y^2+a_z^2\\
			B&=2(a_yV_y+a_zV_z)\\
			C&=V_y^2+V_z^2-R^2\\
			D&=B^2-4AC\\
			t_{1,2}&=\frac{-B\pm\sqrt{D}}{2A}
		\end{split}
	\end{equation}
	
	Подставляем эти $t_{1,2}$ в параметрическое уравнение и получаем 2 точки
	пересечения.
	
	Из них берём ту, которая неотрицательная (не сзади от камеры), минимальную
	(ближе к камере) и соответствует второму неравенству ($C_x\leq{}x\leq{}C_x+H$).
	
	Формулы для остальных осей получаются циклической перестановкой $(x,y,z)$ и
	$(a_x,a_y,a_z)$.
	\newpage
	\section{Пересечение прямой и конуса}
	Рассмотрим конус, направленный вдоль оси $X$, центр основания которого
	расположен в точке $C(C_x;C_y;C_z)$. Радиус - $R$; высота - $H$. Он выражается
	формулами:
	\begin{equation}
		\begin{cases}
			(y-C_y)^2+(z-C_z)^2=(R\frac{C_x+H-x}{H})^2\\
			C_x\leq{}x\leq{}C_x+H
		\end{cases}
	\end{equation}
	
	Назовём начало луча точкой $M(x_0;y_0;z_0)$.
	
	Заранее обозначим $\frac{R^2}{H^2}$ как $U$.
	
	Подставим в это уравнение наши $x$, $y$ и $z$ из параметрического уравнения
	прямой, доведём до квадратного уравнения:
	\begin{equation}
		\begin{split}
			&(y_0+a_yt-C_y)^2+(z_0+a_zt-C_z)^2=U(C_x+H-x_0-a_xt)^2\\\\
			&y_0^2+a_y^2t^2+C_y^2+2y_0a_yt-2y_0C_y-2a_ytC_y\\
			&+z_0^2+a_z^2t^2+C_z^2+2z_0a_zt-2z_0C_z-2a_ztC_z\\
			=&U(C_x^2+H^2+x_0^2+a_x^2t^2+2C_xH\\
			&-2C_xx_0-2C_xa_xt-2Hx_0-2Ha_xt+2x_0a_xt)\\\\
			&t^2(a_y^2+a_z^2-Ua_x^2)\\
			&+t(2y_0a_y-2a_yC_y+2z_0a_z-2a_zC_z+2UC_xa_x+2UHa_x-2Ux_0a_x)\\
			&+1(y_0^2+C_y^2-2y_0C_y+z_0^2+C_z^2-2z_0C_z\\
			&-UC_x^2-UH^2-Ux_0^2-2UC_xH+2UC_xx_0+2UHx_0)=0\\\\
			&t^2(a_y^2+a_z^2-Ua_x^2)\\
			&+t(2(a_y(y_0-C_y)+a_z(z_0-C_z)+Ua_x(C_x+H-x_0)))\\
			&+1((y_0-C_y)^2+(z_0-C_z)^2-U(C_x+H-x_0)^2)=0
		\end{split}
	\end{equation}
	
	Обозначим $C_x+H-x_0$, $y_0-C_y$ и $z_0-C_z$, как $V_x$, $V_y$ и $V_z$
	соответственно, тогда уравнение приобретёт следующий вид:
	\begin{equation}
		\begin{split}
			&t^2(a_y^2+a_z^2-Ua_x^2)+t(2(a_yV_y+a_zV_z+Ua_xV_x))\\
			&+1(V_y^2+V_z^2-UV_x^2)=0
		\end{split}
	\end{equation}
	
	Решаем квадратное уравнение и находим t:
	\begin{equation}
		\begin{split}
			A&=a_y^2+a_z^2-Ua_x^2\\
			B&=2(a_yV_y+a_zV_z+Ua_xV_x)\\
			C&=V_y^2+V_z^2-UV_x^2\\
			D&=B^2-4AC\\
			t_{1,2}&=\frac{-B\pm\sqrt{D}}{2A}
		\end{split}
	\end{equation}
	
	Подставляем эти $t_{1,2}$ в параметрическое уравнение и получаем 2 точки
	пересечения.
	
	Из них берём ту, которая неотрицательная (не сзади от камеры), минимальную
	(ближе к камере) и соответствует второму неравенству ($C_x\leq{}x\leq{}C_x+H$).
	
	Формулы для остальных осей получаются циклической перестановкой $(x,y,z)$ и
	$(a_x,a_y,a_z)$.
\end{document}
